\subsection{A concrete FT215 follow-up}
The new table changes native rows $26/0:0R134\mapsto0R194$ and
$26/1:1R38\mapsto1L114$, leaving all other rows and the source unchanged.
It still has 215 working states. An independent local checker computes 971,847
reachable abstract pairs of a control state and a 15-bit window, beginning with
all windows at state 0. Sliding a window supplies both possible exterior bits,
so this is an inductive overapproximation for every initial tape, including the
blank tape. For all 4,736 eligible row 26 windows, positive old/new traces join
with identical control, head and tape before any exterior read. No additional
modified row is executed in the strict new-trace interior. Exact positive block
matching therefore transports halting in both directions. This is an ordinary
finite mathematical certificate, not a proof-assistant theorem.

For row 26/0, the next bit is 0 and both paths take two steps:
$26\to 134\to17$ and $26\to194\to17$.
For row 26/1, the next two bits are 00; the verified old/new durations are
$(3,5)$ for 1,920 windows and $(5,7)$ for 128 windows.
Fresh symbolic verification independently proves all 6,152 arbitrary-natural
primitive cases on each exact table, using the inherited whole-tape induction
rules. The original/new rule counts are 1,617,650/1,618,486. The exact blank
checkpoint at old step 70,506 is reached at new step 70,526.

The two edits free engine rows 38/0 and 134/0 but introduce obligations at
46/0, 107/1, 114/0, 114/1, 126/0 and 194/0. Thus engine rows increase 144 to 148 and
engine states 80 to 82. This makes the two-step 26/0 shortcut a more attractive
future starting point than the longer 26/1 detour.

A bounded global allocation retains the current source and its 15 contextual
aliases and protects 210 engine, helper and initial-bootstrap rows. Its control
reconstructs exactly the 215-state table. Attempts to delete 38 or 134 both fail.
Independently of the MILP solver, each attempt exhibits 115 distinct decoder
roles whose entire allowed union has only 114 physical slots. The full Hall
witness is included. This proves an obstruction only for the stated pinned
alias and protected row model; it is not a 214-state impossibility result.

After the edits, the remaining entrances are native row $188/1$ to state 38 at source prefix 11,
whose left child is CI6, and native row $312/1$ to state 134 at prefix 001, whose left child is CD3.
Both right children are nonterminal. Native labels are not consecutive: label 312 is a physical state, while the table still has only 215 distinct working labels. Further progress must change the source
role structure, introduce another contextual alias, or separately reimplement
one of the protected engine/helper obligations. The new four-shape rho8 theorem
suggests another emitter language, but its nested implication constructors and
branching must be costed in a literal allocation before claiming a state saving.
